% 超入門：SQUID実験をゼロから理解する
\documentclass[12pt,a4paper]{article}

\usepackage{fontspec}
\usepackage{xeCJK}
\setCJKmainfont{Hiragino Mincho ProN}
\setCJKsansfont{Hiragino Kaku Gothic ProN}
\setCJKmonofont{Hiragino Kaku Gothic ProN}
\usepackage[margin=2.5cm]{geometry}
\usepackage{amsmath,amssymb}
\usepackage{hyperref}
\usepackage{tcolorbox}
\usepackage{enumitem}
\usepackage{booktabs}

\tcbuselibrary{breakable,skins}

\newtcolorbox{naruhodox}[1][]{
  colback=blue!5, colframe=blue!60!black,
  fonttitle=\bfseries\sffamily, title={#1},
  breakable, arc=3mm
}
\newtcolorbox{tatoebabox}[1][]{
  colback=orange!5, colframe=orange!60!black,
  fonttitle=\bfseries\sffamily, title={#1},
  breakable, arc=3mm
}
\newtcolorbox{tukurubox}[1][]{
  colback=green!5, colframe=green!50!black,
  fonttitle=\bfseries\sffamily, title={#1},
  breakable, arc=3mm
}
\newtcolorbox{kanenobox}[1][]{
  colback=yellow!8, colframe=yellow!60!black,
  fonttitle=\bfseries\sffamily, title={#1},
  breakable, arc=3mm
}
\newtcolorbox{abunaibox}[1][]{
  colback=red!5, colframe=red!60!black,
  fonttitle=\bfseries\sffamily, title={#1},
  breakable, arc=3mm
}
\newtcolorbox{honestbox}[1][]{
  colback=gray!8, colframe=gray!50!black,
  fonttitle=\bfseries\sffamily, title={#1},
  breakable, arc=3mm
}

\setlength{\parskip}{0.8em}
\setlength{\parindent}{0pt}

\begin{document}

\begin{center}
{\Huge\bfseries SQUID実験}\\[0.3cm]
{\Huge\bfseries ゼロからわかるガイド}\\[0.8cm]
{\Large 「何もない空間」のエネルギーを測る}\\[0.5cm]
{\Large --- 中学生でも読めるように書きました ---}\\[1.5cm]
{\large Wright Brothers Research Program}\\[0.3cm]
{2026年4月}
\end{center}

\vspace{0.5cm}
\tableofcontents
\newpage

%=====================================================================
\part{そもそも何をしたいのか}
%=====================================================================

\section{「何もない空間」は本当に何もない？}

あなたの目の前にある空間。空気を全部抜いて、光も入れず、
温度もゼロに近づけたとする。完璧な「真空」。

\textbf{直感}: 何もないんだから、エネルギーはゼロ。

\textbf{量子力学の答え}: いいえ、ゼロではありません。

\begin{naruhodox}[なぜ「何もない」のにエネルギーがあるのか]
量子力学には「不確定性原理」というルールがある。

\begin{quote}
\textbf{位置と速度を同時に完全には決められない。}
\end{quote}

これは粒子だけでなく、空間そのものにも当てはまる。
空間の各点で、電磁場が「ちょっとだけ揺れている」。
この揺れは止められない——止めたら不確定性原理に反する。

この「止められない揺れ」がエネルギーを持っている。
これが\textbf{真空エネルギー}（ゼロ点エネルギー）。
\end{naruhodox}

\begin{tatoebabox}[たとえ話：ギターの弦]
ギターの弦を考えよう。弦を弾くと音が出る。
でも弾かなくても、弦は完全に静止はしていない——
空気の分子がぶつかったり、微小な振動が残ったりする。

量子力学の真空は「宇宙サイズのギターの弦」のようなもの。
誰も弾いていないのに、弦は微小に振動し続けている。
その振動のエネルギーが真空エネルギー。

弦の振動には「モード」がある：
基本振動（ドーン）、第2倍音（ドーン×2）、第3倍音（ドーン×3）、\ldots

真空も同じ。モード1、モード2、モード3、\ldots があり、
それぞれがエネルギーを持っている。
\end{tatoebabox}

\section{実験で確認できるのか}

「真空にエネルギーがある」というのは理論の話ではない。
実験で確認されている。それが\textbf{カシミール効果}。

\begin{naruhodox}[カシミール効果]
2枚の金属板を、ものすごく近く（1ミクロン以下）に置く。
すると、2枚の板が\textbf{引き合う力}が生じる。

なぜか？ 板の間には「短い波長のモード」しか入れないから、
板の外より板の間の方が真空エネルギーが低い。
エネルギーが低い方に向かって力が働く
→ 板が引き合う。

この力は1997年にLamoreaux によって精密に測定された。
理論予測と1\%以内で一致。
\textbf{真空エネルギーは実在する。}
\end{naruhodox}

\section{我々は何を測りたいのか}

カシミール効果は「板で一部のモードを制限した」実験。
我々がやりたいのは、もう少し精密なこと：

\begin{center}
\fbox{\parbox{12cm}{\centering\large
\vspace{0.3cm}
\textbf{特定のモード（例えばモード2）だけを消したとき、}\\
\textbf{真空エネルギーはどう変わるか？}
\vspace{0.3cm}
}}
\end{center}

我々の理論の予測：
\begin{itemize}
\item モード2を消すと、真空エネルギーの\textbf{符号が反転}する
（正 → 負に変わる）
\item これは「素数チャンネルのミューティング」という
操作に対応する
\end{itemize}

もし符号が本当に反転したら、それは
「宇宙の真空が整数の数学的構造で決まっている」
ことの最初の実験的証拠になる。

\begin{honestbox}[正直に]
この予測が正しいかどうかはまだわかりません。
だからこそ実験する必要があります。
実験して「変わらなかった」なら、我々の理論は間違いです。
科学とはそういうものです。
\end{honestbox}

%=====================================================================
\part{装置の全体像}
%=====================================================================

\section{何が必要か——5つのパーツ}

この実験装置は、5つのパーツからできている。
一つずつ説明する。

\begin{center}
\renewcommand{\arraystretch}{1.5}
\begin{tabular}{cl}
\toprule
パーツ & 役割 \\
\midrule
1. 箱（キャビティ） & 真空のモードを閉じ込める「部屋」 \\
2. フィルタ（SQUID） & 特定のモードだけを消す「イコライザー」 \\
3. 冷却器 & 全体を $-272$°Cまで冷やす「冷蔵庫」 \\
4. 測定器 & 真空の状態を読み取る「体温計」 \\
5. 配線 & 全部をつなぐ「神経系」 \\
\bottomrule
\end{tabular}
\end{center}

以下、それぞれを日常のたとえで説明する。

%=====================================================================
\part{パーツ1：箱（キャビティ）}
%=====================================================================

\section{キャビティ = 電子レンジの箱}

\begin{tatoebabox}[たとえ話：電子レンジ]
電子レンジの中は金属の箱。マイクロ波（電磁波）が
箱の中で反射を繰り返し、特定の周波数で「共鳴」する。

電子レンジ：大きさ30cm → 共鳴周波数 約2.45 GHz
我々の箱：大きさ3.5cm → 共鳴周波数 約6 GHz

原理は全く同じ。サイズが違うだけ。
\end{tatoebabox}

金属の箱に電磁波（マイクロ波）を入れると、
箱のサイズに合った特定の周波数だけが「きれいに共鳴」する。

\begin{naruhodox}[共鳴とは]
お風呂で「ウォーン」と声を出すと、
特定の高さの声がやたら大きく響くことがある。
これが共鳴。お風呂のサイズに合った波長の音が増幅される。

電磁波でも同じ。箱のサイズに合った波長 → 共鳴 → 強い信号。
合わない波長 → 消える。

我々の箱（内寸 35mm）では：
\begin{itemize}
\item モード1：約 6 GHz（基本共鳴）
\item モード2：約 12 GHz（第2倍音）
\item モード3：約 18 GHz（第3倍音）
\item \ldots
\end{itemize}
\end{naruhodox}

\section{なぜアルミで作るのか}

\begin{naruhodox}[超伝導のすごさ]
普通の金属（銅など）には電気抵抗がある。
電磁波が壁にぶつかるたびにエネルギーが少し失われる。

アルミを $-272$°C以下に冷やすと「超伝導」になる。
電気抵抗がゼロ。
電磁波が壁にぶつかってもエネルギーが失われない。
→ 共鳴がとてつもなく鋭くなる（Q値が100万以上）。

鋭い共鳴 = 各モードをはっきり区別できる
= 「モード2だけを消す」という精密操作が可能になる。
\end{naruhodox}

\begin{tatoebabox}[たとえ話：Q値とは]
Q値 = 「共鳴がどれくらい鋭いか」の数値。

\textbf{Q値が低い}（= 100）：
お風呂の共鳴。「ウォーン」と鳴るが、音を止めるとすぐ消える。
隣の音程との区別がつきにくい。

\textbf{Q値が高い}（= 1,000,000）：
ワイングラスの共鳴。「キーン」と鳴り、指を離しても何秒も鳴り続ける。
他の音程と完全に区別できる。

超伝導キャビティはワイングラス的。
各モードが「キーン」と独立に鳴っている。
だから「モード2だけ消す」という操作ができる。
\end{tatoebabox}

\section{箱の作り方}

\begin{tukurubox}[必要なもの]
\begin{itemize}
\item アルミ合金ブロック（6061-T6）：50mm $\times$ 50mm $\times$ 30mm ← ¥500
\item フライス盤へのアクセス（大学の工作室、またはホームセンター） ← ¥5,000
\item SMAコネクタ $\times$ 2（電磁波の入出力端子） ← ¥2,000
\item インジウムワイヤー（蓋の密封用） ← ¥1,000
\item ネジ（M3 $\times$ 8本） ← ¥200
\end{itemize}
小計：約 ¥8,700
\end{tukurubox}

\begin{tukurubox}[手順]
\begin{enumerate}
\item \textbf{ブロックを2つに分ける}（本体20mm + 蓋10mm）

バンドソーか金鋸で切る。まっすぐ切れなくても大丈夫。

\item \textbf{本体に穴を掘る}（35mm $\times$ 35mm $\times$ 10mm の四角い穴）

フライス盤で削る。精度は $\pm 0.1$mm あれば十分。

フライス盤がない場合：外注。
「アルミブロックに35mm角の穴を深さ10mmで掘ってください」
と図面を渡せば、町工場で3万〜5万円でやってくれる。

\item \textbf{コネクタ穴を開ける}

左右の壁の中央に、直径3mmの穴をドリルで開ける。
ここにSMAコネクタをねじ込む。

\item \textbf{ネジ穴を開ける}

蓋を固定するためのネジ穴（M3）を8箇所。
$\phi$2.5mmドリルで下穴 → M3タップでネジを切る。

\item \textbf{洗浄}

切削油や汚れを落とす。
アセトン（薬局で買える）に10分浸す → IPAで洗う → 水で洗う → 乾かす。

\item \textbf{組み立て}

蓋との接合面にインジウムワイヤーを置いて、蓋をネジで締める。
インジウムは柔らかい金属で、潰れて隙間を塞ぐ。
（ゴムパッキンのようなもの。ただし金属なので超伝導になる。）
\end{enumerate}
\end{tukurubox}

\begin{tatoebabox}[工作のレベル感]
「フライス盤で四角い穴を掘る」は、
大学の工学部1年生が工作実習で最初にやる程度の作業。
精密加工の経験は不要。
ネジ穴あけも、ホームセンターのDIYコーナーで十分可能。

最悪、加工を全部外注しても3〜5万円。
この実験で最も安い部分。
\end{tatoebabox}

%=====================================================================
\part{パーツ2：フィルタ（SQUID）}
%=====================================================================

\section{SQUIDとは何か}

SQUID（スクイッド）= Superconducting QUantum Interference Device
（超伝導量子干渉素子）。

\begin{tatoebabox}[たとえ話：SQUIDは「周波数ノブ付きの消音器」]
カラオケの「ボーカルキャンセル機能」を知っているだろうか。
音楽からボーカルだけを消す機能。

SQUIDはこれの電磁波版。
「モード2（12 GHz）だけを消す」ことができる。
しかも、外から磁場をかけることで
「どの周波数を消すか」をリアルタイムで変えられる。

ツマミ（= 磁場）を回すと、
消す周波数が変わる。
\end{tatoebabox}

\section{原理：なぜ磁場で周波数を変えられるのか}

\begin{naruhodox}[ジョセフソン接合]
2つの超伝導体を、とても薄い（1ナノメートル！）絶縁体で
隔てると、電子のペアが絶縁体を「すり抜ける」。
これをジョセフソン効果という。

普通なら絶縁体は電流を通さない。
でも量子力学の「トンネル効果」により、
超薄い壁なら電子が壁を通り抜ける。

このジョセフソン接合は「磁場に敏感なインダクタ」として働く。
インダクタ = コイルのようなもの。
インダクタの値が変わると、回路の共鳴周波数が変わる。
\end{naruhodox}

\begin{naruhodox}[SQUIDの構造]
SQUID = 「2つのジョセフソン接合をリング状につないだもの」。

リングの中を通る磁場の量（磁束）に応じて、
SQUIDの電気的性質が変わる。

\textbf{磁束 = 0}のとき：SQUID は普通のインダクタ \\
\textbf{磁束 = $\Phi_0/2$}のとき：SQUID のインダクタンスが無限大
→ 電磁波を完全にブロック

$\Phi_0 = 2.07 \times 10^{-15}$ Wb（ウェーバー）= 磁束量子。
とても小さい。地磁気の100億分の1くらい。
だからSQUIDは超高感度の磁気センサーとしても有名。
\end{naruhodox}

\begin{tatoebabox}[たとえ話：磁束量子]
$\Phi_0$ は「超伝導の世界での磁場の最小単位」。

水は一滴ずつしか存在できない（分子があるから）。
同じように、超伝導リングを通る磁場は
$\Phi_0$ の整数倍でしか存在できない。

この「磁場の量子化」が、SQUIDの超高感度の源。
\end{tatoebabox}

\section{SQUIDの作り方——正直に言うと難しい}

\begin{honestbox}[ここが一番難しい]
キャビティ（箱）は工作レベルで作れるが、
SQUIDは\textbf{ナノテク}。
1ナノメートル（髪の毛の10万分の1）の精度が必要。

作り方は2つ：
\begin{enumerate}
\item \textbf{大学のクリーンルームを使う}（推奨）
\item 既製品を買う（高い）
\end{enumerate}
\end{honestbox}

\subsection{方法A：大学のクリーンルームで作る}

\begin{naruhodox}[クリーンルームとは]
空気中のホコリを極限まで少なくした部屋。
半導体工場や大学の理工学部にある。

なぜ必要か？ 1ナノメートルの膜を作るとき、
ホコリ（数マイクロメートル）が1粒あるだけで
接合が壊れる。手術室よりきれいな環境が必要。

多くの大学には「共同利用ナノファブ施設」があり、
学外者でも使用料を払えば使える。
\end{naruhodox}

\begin{tukurubox}[SQUIDの作り方（概要）]
\begin{enumerate}
\item \textbf{シリコン基板を用意}（10mm $\times$ 10mm）

半導体ウエハーを切ったもの。¥3,000/5枚。

\item \textbf{レジスト塗布}

基板表面にPMMA（アクリル樹脂の一種）を薄く塗る。
スピンコーター（回転台）に載せて高速回転させると
均一に薄く（数百ナノメートル）塗れる。

\item \textbf{パターン描画}

電子ビームリソグラフィ（EBL）でSQUIDの形を描く。
光ではなく電子の細いビームで描くので、
数十ナノメートルの精度が出る。

EBLは大学の共同利用施設で使える（¥20,000/回くらい）。

パターンのデザインは無料ソフト（KLayout）で作る。

\item \textbf{現像}

描いた部分のレジストを溶かす。
MIBK:IPA = 1:3 の液に1分浸すだけ。
お菓子の型抜きのようなもの。

\item \textbf{アルミ蒸着（ここが核心）}

真空チャンバーの中で、アルミを加熱して蒸発させ、
基板に薄い膜（30ナノメートル）を付ける。

\textbf{ポイント}：斜めの角度から蒸着する。
\begin{itemize}
\item 1回目：右斜め30°からアルミを30nm蒸着
\item 酸素を入れて表面を酸化（1ナノメートルの絶縁膜ができる）
\item 2回目：左斜め30°からアルミを50nm蒸着
\end{itemize}

2回の蒸着が「ブリッジ」（レジストの残った部分）の下で
重なる場所がジョセフソン接合になる。
間に挟まった酸化アルミ（1nm）が絶縁体。

\item \textbf{リフトオフ}

レジストを溶かして（NMPという溶剤に浸す）、
レジストの上に乗ったアルミごと除去する。
基板上に直接付いたアルミだけが残る。

これでSQUIDの完成。
\end{enumerate}
\end{tukurubox}

\begin{abunaibox}[最も難しいポイント：酸化条件]
ジョセフソン接合の性能は「酸化膜の厚さ」で決まる。

\begin{itemize}
\item 酸化が足りない → 膜が薄すぎ → 電流が流れすぎ → 壊れる
\item 酸化しすぎ → 膜が厚すぎ → 電流が流れない → 動かない
\item ちょうどいい → 臨界電流 $I_c \approx 1 \mu$A → 完璧
\end{itemize}

\textbf{対策}：テストピースを5枚作り、酸化条件を変えて（酸素圧を
100/200/300/400/500 mTorr）、室温で抵抗値を測る。
$R \approx 440 \Omega$ になる条件が当たり。

失敗しても基板代は¥600/枚。5回やっても¥3,000。
何度でもやり直せる。
\end{abunaibox}

\subsection{方法B：既製品を買う}

\begin{kanenobox}[市販SQUIDチップ]
Star Cryoelectronics（アメリカ）などが
SQUIDチップを販売している。

価格：1個 約5万〜15万円。

メリット：確実に動く。\\
デメリット：高い。設計の自由度が低い。

\textbf{妥協案}：最初の実験は市販品で「装置全体が動くか」を確認し、
2回目以降で自作SQUIDに切り替える。
\end{kanenobox}

%=====================================================================
\part{パーツ3：冷却器}
%=====================================================================

\section{なぜ冷やす必要があるのか}

\begin{tatoebabox}[たとえ話：星空の観察]
星空を見たいとする。

\textbf{昼間}：太陽が明るすぎて星が見えない。\\
\textbf{夜}：太陽がなくなって、星の光だけが見える。

我々の実験でも同じ。

\textbf{温度が高い}：「熱の光」が強すぎて、
真空のかすかなエネルギーが見えない。\\
\textbf{温度を下げる}：「熱の光」が消えて、
真空エネルギーだけが見える。
\end{tatoebabox}

理由は2つ：

\textbf{理由1：超伝導のため。}
アルミが超伝導になるのは $-272$°C（1.2 K）以下。
冷やさないと箱が超伝導にならず、Q値が低いままで
モードを区別できない。

\textbf{理由2：真空を「見る」ため。}
温度が高いと「熱的な光子」（熱の粒）がたくさんいて、
真空の揺らぎが隠れる。
$-272.7$°C（0.3 K）まで冷やすと、
熱的光子がほぼゼロになり、真空だけが残る。

\section{冷却器の種類}

\begin{naruhodox}[冷やし方の階層]
\begin{center}
\renewcommand{\arraystretch}{1.4}
\begin{tabular}{lll}
\toprule
方法 & 温度 & 日常のたとえ \\
\midrule
冷蔵庫 & $+5$°C & 普通の冷蔵庫 \\
冷凍庫 & $-20$°C & アイスを保存 \\
ドライアイス & $-78$°C & 宅配のアイス \\
液体窒素 & $-196$°C & テレビの実験でバラを凍らせるやつ \\
液体ヘリウム & $-269$°C & ★ ここからが極低温の世界 \\
$^3$He冷凍機 & $-272.7$°C & ★★ 我々が使うのはこれ \\
希釈冷凍機 & $-272.99$°C & ★★★ Google の量子コンピュータ \\
\bottomrule
\end{tabular}
\end{center}

「$-272.7$°C」と「$-273$°C」（絶対零度）の差はわずか0.3度。
絶対零度のたった0.3度上まで冷やす。
\end{naruhodox}

\begin{naruhodox}[$^3$He冷凍機の仕組み]

ヘリウムには2種類ある：
\begin{itemize}
\item $^4$He：普通のヘリウム。沸点 $-269$°C（4.2 K）。風船に入ってるやつ。
\item $^3$He：軽いヘリウム（中性子が1個少ない）。沸点がもっと低い。
\end{itemize}

$^3$He冷凍機の原理：
\begin{enumerate}
\item まず液体 $^4$He で $-269$°C まで冷やす
\item 液体 $^3$He をポンプで吸い出す（蒸発冷却）
\item 蒸発するときに周りから熱を奪う（汗が蒸発して涼しくなるのと同じ）
\item $-272.7$°C（0.3 K）に到達
\end{enumerate}

ポイント：$^3$He は非常に希少（天然にほとんどない）。
数リットルの $^3$He がシステム内を循環する。
高価だが、\textbf{気体なので漏らさなければ永久に使える}。
\end{naruhodox}

\begin{kanenobox}[冷却器のコスト]
$^3$He冷凍機を\textbf{持っている}大学の共同利用施設を使う。

使用料の目安：
\begin{itemize}
\item 1回のクールダウン（冷やして→測定して→戻す）：3〜10万円
\item 液体ヘリウム代：約50リットル必要、¥30,000〜¥50,000
\item 5回の実験で：約30万円
\end{itemize}

自分で冷凍機を買うと：1000万〜3000万円。買ってはいけない。

\textbf{使える大学の例}：
東大物性研、東北大金研、阪大極低温センター、
筑波大、京大、名大、東工大、北大など。
多くの大学に低温実験施設がある。
\end{kanenobox}

%=====================================================================
\part{パーツ4：測定器}
%=====================================================================

\section{何を測るのか}

\begin{naruhodox}[S$_{21}$ パラメータ]
箱の入口からマイクロ波を入れて、出口から出てくる量を測る。
入力に対する出力の比を $S_{21}$ と呼ぶ。

共鳴周波数では：$S_{21}$ が大きくなる（波が通りやすい）。
共鳴以外では：$S_{21}$ が小さい（波が通りにくい）。

$S_{21}$ を周波数の関数として描くと、
共鳴のところに「ピーク」が立つ。
このピークの位置、高さ、幅から、
キャビティの中で何が起きているかがわかる。
\end{naruhodox}

\begin{naruhodox}[VNA（ベクトルネットワークアナライザー）]
VNA = $S_{21}$ を測る装置。
マイクロ波を出し、受け取り、振幅と位相を記録する。

周波数範囲：100 MHz 〜 20 GHz くらい。
我々の実験では 4 〜 14 GHz を使う。

VNA は高価（新品 500万〜2000万円）だが、
大学の共同利用なら使用料 ¥50,000 程度。
\end{naruhodox}

\section{測定の流れ}

\begin{tukurubox}[ステップバイステップ]
\begin{enumerate}
\item \textbf{装置を冷やす}（朝〜昼）

冷凍機を起動して $-272.7$°C に到達するまで待つ。
6〜8時間。

\item \textbf{共鳴を確認する}

VNAで 4〜14 GHz をスキャン。
$S_{21}$ のグラフにピーク（共鳴）が見えるはず。

\begin{itemize}
\item モード1：約 6 GHz
\item モード2：約 12 GHz
\end{itemize}

ピークが見えたら、装置は正常に動いている。

\item \textbf{SQUIDをチューニング}

磁束バイアスコイル（SQUIDの周りに巻いた銅線）に
電流を流す。電流を変えると、SQUIDの共鳴周波数が変わる。

SQUIDの周波数がモード2（12 GHz）に一致するまで電流を調整。
→ $S_{21}$ のモード2のピークが「消える」（SQUIDが吸収）。

\item \textbf{測定A：SQUIDオフ}

SQUIDをモード2から外す（電流を変えて離調）。
→ 全モード生きている状態。
→ $S_{21}$ を記録。

\item \textbf{測定B：SQUIDオン}

SQUIDをモード2に合わせる。
→ モード2だけ消えた状態。
→ $S_{21}$ を記録。

\item \textbf{AとBを比較}

モード1（6 GHz）のピーク位置は、AとBで違うか？
もし違えば：「モード2を消したことで、モード1の真空エネルギーが変わった」。
\end{enumerate}
\end{tukurubox}

\begin{tatoebabox}[たとえ話：オーケストラのチューニング]
オーケストラで、オーボエがA（440 Hz）を鳴らし、
他の楽器がそれに合わせる。

もし\textbf{バイオリンだけ退場させたら}、
オーボエのAの音は変わるだろうか？

普通は変わらない（各楽器は独立だから）。
でも、もし楽器同士が「真空を通じて」つながっていたら？
→ バイオリンを消すことで、オーボエの音が微妙に変わるかもしれない。

これが我々の実験。「モード2を消すとモード1が変わるか」。
\end{tatoebabox}

%=====================================================================
\part{パーツ5：配線}
%=====================================================================

\section{なぜ特殊なケーブルが必要か}

\begin{naruhodox}[低温配線の問題]
普通のケーブル（銅の同軸ケーブル）は熱をよく伝える。
室温（$+20$°C）から $-272.7$°C の装置に銅ケーブルをつなぐと、
熱が流れ込んで冷凍機が温まってしまう。

対策：
\begin{itemize}
\item \textbf{ステンレス同軸ケーブル}：銅より1000倍熱を伝えにくい
\item \textbf{アッテネータ}（減衰器）：各温度段に挿入。
信号を弱めるが、同時に熱ノイズも減衰させる
\item \textbf{超伝導ケーブル}：最低温部分にはNbTi合金の
超伝導ケーブルを使う。抵抗ゼロ＋低熱伝導
\end{itemize}
\end{naruhodox}

\begin{kanenobox}[ケーブルのコスト]
低温用同軸ケーブル $\times$ 2本：約 ¥30,000

普通のケーブルの10倍するが、これがないと実験できない。
\end{kanenobox}

%=====================================================================
\part{全体のコスト}
%=====================================================================

\section{最小構成：約57万円}

\begin{kanenobox}[費用内訳]
\begin{center}
\renewcommand{\arraystretch}{1.4}
\begin{tabular}{lr}
\toprule
項目 & 金額 \\
\midrule
キャビティ（箱）材料 + 加工 & ¥9,000 \\
SQUIDチップ（自作の場合） & ¥48,000 \\
\quad（市販品を買う場合） & \quad（¥100,000） \\
磁束バイアス部品 & ¥500 \\
$^3$He冷凍機使用料（5回） & ¥300,000 \\
液体ヘリウム & ¥50,000 \\
VNA使用料 & ¥50,000 \\
低温用同軸ケーブル & ¥30,000 \\
その他消耗品 & ¥10,000 \\
\midrule
\textbf{合計} & \textbf{¥497,500} \\
\bottomrule
\end{tabular}
\end{center}

端数を切り上げて\textbf{約50万〜60万円}。

市販SQUIDを使う場合：+5万で約55万〜65万円。

\textbf{比較}：
\begin{itemize}
\item CERNのLHC：約6000億円
\item JamesWebb宇宙望遠鏡：約1兆円
\item この実験：約60万円（自家用車より安い）
\end{itemize}
\end{kanenobox}

\begin{honestbox}[隠れコストに注意]
上の表に含まれて\textbf{いない}もの：
\begin{itemize}
\item あなたの人件費（何百時間もかかる）
\item 失敗したときの材料費（SQUID作りは10回失敗する覚悟で）
\item 交通費（施設が遠い場合）
\item 大学施設への登録費用（場所による）
\end{itemize}
現実的には、予備費込みで \textbf{80万〜120万円} を見ておく方が安全。
\end{honestbox}

%=====================================================================
\part{結果の見方}
%=====================================================================

\section{何が起きたら「成功」か}

\begin{tukurubox}[Level 0：装置が動いた]
SQUIDのオン/オフで $S_{21}$ が変化する。

意味：装置は正常。SQUIDはフィルタとして機能している。
新しい物理ではないが、第一歩。
\end{tukurubox}

\begin{tukurubox}[Level 1：真空エネルギーが変わった ★]
モード2をSQUIDで消したとき、
モード1の共鳴周波数が\textbf{統計的に有意に}シフトした。

意味：「1つのモードを消すと真空のエネルギーが変わる」。
これは\textbf{算術的真空工学の最初の実験的証拠}。
論文になる。ニュースになる。

予想されるシフト量：数 kHz 〜 数十 kHz（検出可能）。
\end{tukurubox}

\begin{tukurubox}[Level 2：符号が反転した ★★★]
モード2を消したとき、モード1のシフトの\textbf{方向}が
理論予測（正→負）と一致。

意味：「宇宙の真空は整数の構造で記述される」という仮説の
強力な証拠。物理学のパラダイムシフト。
ノーベル賞級（もしかしたら）。
\end{tukurubox}

\begin{honestbox}[何が起きたら「失敗」か]
何も変わらなかった場合、
我々の理論（少なくとも実験に関する予測）は間違い。

しかし「何も起きなかった」という結果にも価値がある。
「この方法では真空エネルギーの変化は検出できない」
という事実は、将来の研究者にとって重要な情報。

科学において、「やってみたら違った」は失敗ではない。
「やらなかった」が唯一の失敗。
\end{honestbox}

%=====================================================================
\part{FAQ（よくある質問）}
%=====================================================================

\section*{Q1：高校生でもできますか？}

キャビティ（箱）の製作は高校生でもできます。
SQUID製作にはクリーンルームへのアクセスが必要で、
通常は大学の施設を使います。
大学の先生にコンタクトして共同研究として進めるのが現実的です。

\section*{Q2：危険はありますか？}

\begin{abunaibox}[安全上の注意]
\begin{itemize}
\item \textbf{液体ヘリウム}：$-269$°C。皮膚に触れると凍傷。
手袋とゴーグル必須。必ず経験者の指導の下で扱うこと。
\item \textbf{液体窒素}：同上。
\item \textbf{真空装置}：内破の危険。ガラス容器は使わない。
\item \textbf{マイクロ波}：人体への影響は微小（出力が非常に小さいため）
だが、直接目に当てないこと。
\item \textbf{薬品}（アセトン、NMP等）：換気の良い場所で使用。
\end{itemize}
\end{abunaibox}

\section*{Q3：一人でできますか？}

理論的には可能ですが、推奨しません。
低温実験は「冷凍機を知っている人」の助けが不可欠です。
大学の物性物理や量子エレクトロニクスの研究室に
コンタクトしてください。

\section*{Q4：どれくらい時間がかかりますか？}

\begin{itemize}
\item キャビティ製作：1〜2日
\item SQUID製作：1〜2週間（慣れれば3日）
\item 冷凍機の予約待ち：1〜3ヶ月
\item 1回の測定（クールダウン〜データ取得）：1〜2日
\item データ解析：1〜2週間
\end{itemize}

準備開始から結果が出るまで：\textbf{3〜6ヶ月}が目安。

\section*{Q5：この実験で宇宙が壊れることはありますか？}

ありません。理論的に証明済みです。

局所化 $\mathbb{Z} \to \mathbb{Z}[1/p]$ による真空の変化は
$K_1$ のトポロジカル不変量で保護されており、
連続的に「普通の真空」に崩壊できません。
つまり、実験室の中で局所的に真空を変えても、
それが宇宙全体に広がることはありません。

それ以前に、この実験のエネルギースケール（$\sim 10^{-24}$ J）は、
宇宙線が日常的に地球に降り注いでいるエネルギー
（$\sim 10^{-7}$ J）の$10^{17}$倍も小さい。
宇宙線でも宇宙は壊れていないので、この実験は全く安全です。

\vfill
\begin{center}
\rule{8cm}{0.4pt}\\[0.5cm]
{\large\textbf{やってみなければ、わからない。}}\\[0.3cm]
{\small Wright Brothers Research Program, 2026}
\end{center}

\end{document}
